\chapter{Leçon 17}

\section{Strukturéiert Übungen (Level A1 / Niveau A1)}

Dës Übunge vum Jérôme Lulling hëllefen den Alldagsvokabulaire an d'Grammatik ze üben. Wielt déi richteg Äntwert fir d'Lück am Saz (A, B, C oder D) auszefëllen. \\
These exercises from Jérôme Lulling help practice everyday vocabulary and grammar. Choose the correct answer to fill the gap in the sentence (A, B, C, or D). \\
Ces exercices de Jérôme Lulling aident à pratiquer le vocabulaire quotidien et la grammaire. Choisissez la bonne réponse pour remplir la lacune dans la phrase (A, B, C ou D).

\begin{enumerate}
    \item Haut ass den 22.04. \\
    \textbf{A)} zweeanzwanzegsten \quad \textbf{B)} aachtandréssegsten \quad \textbf{C)} fënneften \quad \textbf{D)} Abrëll.
    
    \item Haut ass Méinden. Muer ass \_\_\_\_\_\_\_. \\
    \textbf{A)} Mëttwoch \quad \textbf{B)} Donneschden \quad \textbf{C)} Samschden \quad \textbf{D)} Dënschden.
    
    \item Haut ass Méinden. Gëschter war \_\_\_\_\_\_\_. \\
    \textbf{A)} Samschden \quad \textbf{B)} Dënschden \quad \textbf{C)} Sonnden \quad \textbf{D)} Mëttwoch.
    
    \item Et ass elo 12:05. \\
    \textbf{A)} fënnef vir \quad \textbf{B)} zéng vir \quad \textbf{C)} fënnef op \quad \textbf{D)} zwielef.
    
    \item Et ass elo 13:15. \\
    \textbf{A)} fënnef op \quad \textbf{B)} Véierel op \quad \textbf{C)} zéng vir \quad \textbf{D)} eng.
    
    \item Haut ass Samschden. Iwwermuer ass \_\_\_\_\_\_\_. \\
    \textbf{A)} Sonnden \quad \textbf{B)} Mëttwoch \quad \textbf{C)} Méinden \quad \textbf{D)} Dënschden.
    
    \item Haut ass Méinden. Virugëschter war \_\_\_\_\_\_\_. \\
    \textbf{A)} Mëttwoch \quad \textbf{B)} Samschden \quad \textbf{C)} Dënschden \quad \textbf{D)} Donneschden.
    
    \item Haut ass de 15.05. \\
    \textbf{A)} fënnefte \quad \textbf{B)} fofzéngte \quad \textbf{C)} zweete \quad \textbf{D)} Mee.
    
    \item Et ass elo 16:30. \\
    \textbf{A)} hallwer dräi \quad \textbf{B)} hallwer sechs \quad \textbf{C)} hallwer fënnef.
    
    \item Haut ass den 1.04 (Abrëll). \\
    \textbf{A)} zweeten \quad \textbf{B)} drëtten \quad \textbf{C)} éischten.
    
    \item Ech hu \_\_\_\_\_\_\_ Problem. \\
    \textbf{A)} keng \quad \textbf{B)} net \quad \textbf{C)} kee.
    
    \item Wou \_\_\_\_\_\_\_ du? \\
    \textbf{A)} wunnen \quad \textbf{B)} wunns \quad \textbf{C)} wunnt.
    
    \item Ech sinn \_\_\_\_\_\_\_ hongereg. \\
    \textbf{A)} keen \quad \textbf{B)} keng \quad \textbf{C)} net.
    
    \item Ech gesinn dräi \_\_\_\_\_\_\_. \\
    \textbf{A)} Bamen \quad \textbf{B)} Beemer \quad \textbf{C)} Beem \quad \textbf{D)} Beemen.
    
    \item Vu wou \_\_\_\_\_\_\_ du? \\
    \textbf{A)} kënnt \quad \textbf{B)} kënns \quad \textbf{C)} këmms.
    
    \item Wéini \_\_\_\_\_\_\_ du doheem? \\
    \textbf{A)} bass \quad \textbf{B)} buss \quad \textbf{C)} biss.
    
    \item Ech hu \_\_\_\_\_\_\_ Zäit. \\
    \textbf{A)} keen \quad \textbf{B)} keng \quad \textbf{C)} net.
    
    \item Hien huet dräi \_\_\_\_\_\_\_. \\
    \textbf{A)} Autos \quad \textbf{B)} Autoer \quad \textbf{C)} Autoen \quad \textbf{D)} Autoën.
    
    \item Wéini \_\_\_\_\_\_\_ du dat? \\
    \textbf{A)} maachs \quad \textbf{B)} méchs \quad \textbf{C)} maachen.
    
    \item Ech sangen dräi \_\_\_\_\_\_\_. \\
    \textbf{A)} Lidden \quad \textbf{B)} Lidds \quad \textbf{C)} Lidder \quad \textbf{D)} Lidda.
    
    \item Wéi vill Kanner \_\_\_\_\_\_\_ du? \\
    \textbf{A)} huet \quad \textbf{B)} hunns \quad \textbf{C)} hues.
    
    \item Wéi al \_\_\_\_\_\_\_ du? \\
    \textbf{A)} buss \quad \textbf{B)} bass \quad \textbf{C)} biss.
    
    \item Wéi vill Suen \_\_\_\_\_\_\_ hatt? \\
    \textbf{A)} hues \quad \textbf{B)} hutt \quad \textbf{C)} huet.
    
    \item Wat \_\_\_\_\_\_\_ du? \\
    \textbf{A)} sicht \quad \textbf{B)} sichs \quad \textbf{C)} sichen.
    
    \item Ech hu \_\_\_\_\_\_\_ Brëll. \\
    \textbf{A)} keng \quad \textbf{B)} net \quad \textbf{C)} kee.
    
    \item Ech iesse léiwer \_\_\_\_\_\_\_ wéi Orangen. \\
    \textbf{A)} Autoen \quad \textbf{B)} Äppel \quad \textbf{C)} Haiser.
    
    \item Wat \_\_\_\_\_\_\_ du gär? \\
    \textbf{A)} ësst \quad \textbf{B)} iessen \quad \textbf{C)} ëss.
    
    \item Wat \_\_\_\_\_\_\_ du gär? \\
    \textbf{A)} drénkt \quad \textbf{B)} drénks \quad \textbf{C)} drénken.
    
    \item Wéi \_\_\_\_\_\_\_ däin Hond? \\
    \textbf{A)} heescht \quad \textbf{B)} heeschs \quad \textbf{C)} heeschen.
    
    \item Wéini \_\_\_\_\_\_\_ du mech un? \\
    \textbf{A)} rifft \quad \textbf{B)} ruff \quad \textbf{C)} riffs.
    
    \item Ech \_\_\_\_\_\_\_ keng Zäit. \\
    \textbf{A)} hu \quad \textbf{B)} hunn \quad \textbf{C)} hun.
    
    \item Ech schwatzen \_\_\_\_\_\_\_. \\
    \textbf{A)} Heizung \quad \textbf{B)} Däitsch \quad \textbf{C)} Zeitung.
    
    \item Ech si \_\_\_\_\_\_\_. \\
    \textbf{A)} Blumm \quad \textbf{B)} Buedzëmmer \quad \textbf{C)} Kleeder \quad \textbf{D)} räich.
    
    \item Mäi Beruff ass \_\_\_\_\_\_\_. \\
    \textbf{A)} Stull \quad \textbf{B)} Metzler \quad \textbf{C)} giel \quad \textbf{D)} rout.
    
    \item Ech hunn dräi \_\_\_\_\_\_\_. \\
    \textbf{A)} Kand \quad \textbf{B)} aarm \quad \textbf{C)} Kanner.
    
    \item Mat \_\_\_\_\_\_\_ schwätz du? \\
    \textbf{A)} weem \quad \textbf{B)} Auto \quad \textbf{C)} Heizung.
    
    \item Wéini \_\_\_\_\_\_\_ Zäit? \\
    \textbf{A)} huet du \quad \textbf{B)} hues du \quad \textbf{C)} hunn du.
    
    \item Hues du \_\_\_\_\_\_\_ Fra? \\
    \textbf{A)} en \quad \textbf{B)} eng \quad \textbf{C)} e.
    
    \item Hues du \_\_\_\_\_\_\_ Auto? \\
    \textbf{A)} eng \quad \textbf{B)} en \quad \textbf{C)} e.
    
    \item Ech gesinn dräi \_\_\_\_\_\_\_. \\
    \textbf{A)} Kou \quad \textbf{B)} Kouen \quad \textbf{C)} Kéi.
    
    \item Kenns du \_\_\_\_\_\_\_ Tom? \\
    \textbf{A)} d' \quad \textbf{B)} de \quad \textbf{C)} den.
    
    \item Kenns du \_\_\_\_\_\_\_ Chantal? \\
    \textbf{A)} den \quad \textbf{B)} de \quad \textbf{C)} d'.
\end{enumerate}


\section{Iwwersetzung an Erklärungen (Translation and Explanations)}

Hei sinn e puer Iwwersetzungen an Erklärunge fir dat virausgaangent Übungsblat. \\
Here are some translations and explanations for the previous exercise sheet. \\
Voici quelques traductions et explications pour la feuille d'exercices précédente.

\begin{itemize}
    \item \textbf{Ech sinn net hongereg / Ech hu keen Honger} \\
    Fr: Je n'ai pas faim (\textit{Je ne suis pas affamé}) \\
    En: I am not hungry.
    
    \item \textbf{Vu wou kënns du?} \textit{(Opgepasst, et ass mat K net mat M, d'Usprooch ass 'kënns/kënnt')} \\
    Fr: D'où viens-tu ? \\
    En: Where do you come from?
    
    \item \textbf{Wéini bass du doheem?} \\
    Fr: Quand es-tu à la maison ? \\
    En: When are you home?
    
    \item \textbf{Ech iesse léiwer Äppel wéi Orangen.} \\
    Fr: Je préfère manger des pommes plutôt que des oranges. \\
    En: I prefer eating apples over oranges.
    
    \item \textbf{Wéini riffs du mech un?} \textit{(Vum Verb : uruffen)} \\
    Fr: Quand m'appelles-tu ? \\
    En: When are you going to call me?
    
    \item \textbf{Mat weem schwätz du?} \\
    Fr: Avec qui parles-tu ? \\
    En: Who are you talking to?
\end{itemize}

\subsection*{Wichteg Onreegelméisseg Pluralformen (Important Irregular Plurals / Pluriels Irréguliers Importants)}
\begin{itemize}
    \item \textbf{e Bam} (a tree / un arbre) $\rightarrow$ \textbf{dräi Beem} (three trees / trois arbres) 
    \item \textbf{eng Kou} (a cow / une vache) $\rightarrow$ \textbf{zwou Kéi} (two cows / deux vaches)
    \item \textbf{e Kand} (a child / un enfant) $\rightarrow$ \textbf{vill Kanner} (many children / beaucoup d'enfants)
    \item \textbf{en Auto} (a car / une voiture) $\rightarrow$ \textbf{zwee Autoen} (two cars / deux voitures)
\end{itemize}


\section{D'Eifeler Regel (N-Rule / Règle de l'Eifel)}

A very important pronunciation and spelling rule in Luxembourgish is the \textbf{Eifeler Regel}. It dictates when you drop the final \textbf{-n} or \textbf{-nn} of a word (mostly adjectives, articles, and pronouns) depending on the first letter of the following word. \\
\textit{Une règle de prononciation et d'orthographe très importante en luxembourgeois est la \textbf{Eifeler Regel}. Elle dicte quand il faut laisser tomber le \textbf{-n} ou \textbf{-nn} final d'un mot (la plupart du temps des adjectifs, des articles et des pronoms) selon la première lettre du mot suivant.}

\subsection*{Wéini bleift den '-n'? (When does the '-n' stay? / Quand le '-n' reste-t-il ?)}
You \textbf{keep} the final '-n' if the following word starts with a vowel or the specific consonants: \textbf{n, d, t, z, h, a, e, i, o, u}. \\
\textit{Vous \textbf{gardez} le '-n' final si le mot suivant commence par une voyelle ou les consonnes spécifiques : \textbf{n, d, t, z, h, a, e, i, o, u}.}

\textbf{Mnemonic to remember / Moyen mnémotechnique à retenir :} \textbf{"united zoha"} or \textbf{"H U N D E T R A P E Z"} (if you include vowels as part of the sounds / si vous incluez les voyelles). An easy way to recall the consonants is / Un moyen facile de se rappeler des consonnes est : \textbf{n, d, t, z, h} + any vowel / n'importe quelle voyelle.

\textit{Beispiller (Examples / Exemples):}
\begin{itemize}
    \item de\textbf{n} \textbf{H}ond (the dog / le chien)
    \item de\textbf{n} \textbf{D}ësch (the table / la table)
    \item eise\textbf{n} \textbf{A}uto (our car / notre voiture)
    \item e gudde\textbf{n} \textbf{N}oper (a good neighbor / un bon voisin)
\end{itemize}

\subsection*{Wéini fält den '-n' ewech? (When do you drop the '-n'? / Quand laisse-t-on tomber le '-n' ?)}
You \textbf{drop} the final '-n' before all other consonants (b, c, f, g, k, l, m, p, q, r, s, v, w, x, y). \\
\textit{Vous \textbf{laissez tomber} le '-n' final devant toutes les autres consonnes (b, c, f, g, k, l, m, p, q, r, s, v, w, x, y).}

\textit{Beispiller (Examples / Exemples):}
\begin{itemize}
    \item de \textbf{M}ann (the man / l'homme - instead of / au lieu de \textit{den} Mann)
    \item eise \textbf{B}ouf (our boy / notre garçon - instead of / au lieu de \textit{eisen} Bouf)
    \item e gudde \textbf{K}olleg (a good friend / un bon ami - instead of / au lieu de \textit{gudden})
\end{itemize}


\section{Besëtzpronomen (Pronoms Possessifs / Possessive Pronouns)}

Hei sinn d'Besëtzpronomen fir all d'Persounen op Lëtzebuergesch. \\
Here are the possessive pronouns for all persons in Luxembourgish. \\
Voici les pronoms possessifs pour toutes les personnes en luxembourgeois.

\subsection*{Pronomentabell (Pronoun Table / Tableau des pronoms)}

\begin{table}[h!]
    \centering
    \renewcommand{\arraystretch}{1.5}
    \begin{tabular}{|l|c|c|c|c|}
    \hline
    \textbf{Subject (Sujet)} & \textbf{Masc. (Männech)} & \textbf{Fem. (Weiblech)} & \textbf{Neut.} & \textbf{Plural} \\
    \hline
    \textbf{ech} (I / je) & mäin & meng & mäin & meng \\
    \hline
    \textbf{du} (you, sg. / tu) & däin & deng & däin & deng \\
    \hline
    \textbf{hien} (he / il) & säin & seng & säin & seng \\
    \hline
    \textbf{hatt / Si} (she / elle) \textsuperscript{1} & säin / hiren \textsuperscript{2} & seng / hir & säin / hiert & seng / hir \\
    \hline
    \textbf{et} (it / neutre) & säin & seng & säin & seng \\
    \hline
    \textbf{mir} (we / nous) & eisen & eis & eist & eis \\
    \hline
    \textbf{dir} (you, pl. / vous) \textsuperscript{3} & ären & är & äert & är \\
    \hline
    \textbf{si} (they / ils, elles) \textsuperscript{4} & hiren & hir & hiert & hir \\
    \hline
    \textbf{Dir} (You, formal sg/pl / vous formel) & Ären & Är & Äert & Är \\
    \hline
    \textbf{Si} (They, formal / ils, elles formel) & Hiren & Hir & Hiert & Hir \\
    \hline
    \end{tabular}
\end{table}

\textbf{Notten / Notes:}
\begin{enumerate}
    \item \textbf{hatt vs Si (she / elle):} In Luxembourgish, you use "hatt" if you call the person by their first name. Otherwise, use "Si" (capitalized). \\
    \textit{En luxembourgeois, vous utilisez "hatt" si vous appelez la personne par son prénom. Sinon, utilisez "Si" (avec majuscule).}
    \item \textbf{säin vs Hiren/hiren:} If you use "hatt", the possessive pronouns act like the masculine/neuter forms (\textit{säin, seng, säin, seng}).  \\
    \textit{Si vous utilisez "hatt", les pronoms possessifs se comportent comme les formes masculines/neutres.}
    \item The lowercase \textbf{dir} vs the capitalized \textbf{Dir}: The lowercase \textbf{dir} refers to plural you ("y'all" / "vous" au pluriel informel), whereas \textbf{Dir} forms point to the polite "You" ("vous de politesse"). Their pronouns follow the same pattern \textit{ären vs Ären}.
    \item The lowercase \textbf{si} vs capitalized \textbf{Si}: Same as above. \textbf{si} is "they" (ils/elles), whereas \textbf{Si} is polite formal "They" sometimes used historically. Their possessives are \textit{hiren vs Hiren}.
\end{enumerate}

\subsection*{Beispiller (Exemples / Examples)}

\begin{itemize}
    \item \textbf{ech / I / je:} \\
    Mäi Brëll ass net hei. \textit{(My glasses are not here. / Mes lunettes ne sont pas là.)} *Note the Eifeler-Regel drops the '-n' before consonants! / \textit{Remarquez que la loi Eifeler-Regel fait tomber le '-n' devant les consonnes !}* \\
    Ech froe meng Mamm. \textit{(I ask my mother. / Je demande à ma mère.)}
    
    \item \textbf{du / you (sg) / tu:} \\
    Wou ass däin Auto? \textit{(Where is your car? / Où est ta voiture ?)} \\
    Sinn dat deng Kanner? \textit{(Are those your children? / Sont-ce tes enfants ?)}
    
    \item \textbf{hien / he / il:} \\
    Dat ass säi Frënd. \textit{(That is his friend. / C'est son ami.)} \\
    Hie spillt mat senger Schwëster. \textit{(He plays with his sister. / Il joue avec sa sœur.)}
    
    \item \textbf{hatt / Si (she / elle):} \\
    Hatt hëlt säi Buch. \textit{(She takes her book. / Elle prend son livre.) - Familiar / Familier} \\
    Si hëlt hiert Buch. \textit{(She takes her book. / Elle prend son livre.) - Formal / Formel} \\
    Huet hatt seng Hausaufgabe gemaach? \textit{(Did she do her homework? / A-t-elle fait ses devoirs ?)}
    
    \item \textbf{et / it / neutre:} \\
    D'Kand drénkt säi Waasser. \textit{(The child drinks its water. / L'enfant boit son eau.)}
    
    \item \textbf{mir / we / nous:} \\
    Wou ass eisen Dësch? \textit{(Where is our table? / Où est notre table ?)} \\
    Mir ginn an eis Vakanz. \textit{(We are going on our vacation. / Nous partons en vacances.)}
    
    \item \textbf{dir / you (pl) / vous:} \\
    Hutt dir äre Schlëssel? \textit{(Do you have your key? / Avez-vous votre clé ?)} \\
    Ass dat äert Haus? \textit{(Is that your house? / Est-ce votre maison ?)}
    
    \item \textbf{si / they / ils, elles:} \\
    Si wunnen an hirem neien Haus. \textit{(They live in their new house. / Ils habitent dans leur nouvelle maison.)} \\
    Dat sinn hir Autoen. \textit{(Those are their cars. / Ce sont leurs voitures.)}
    
    \item \textbf{Dir / You (formal) / vous formel:} \\
    Ass dat Ären Auto? \textit{(Is this Your car? / Est-ce Votre voiture ?)} \\
    Wou sinn Är Kollegen? \textit{(Where are Your colleagues? / Où sont Vos collègues ?)}
\end{itemize}
